\phantomsection
\section*{修订历史}\label{revision-history}
\lhead{\it{修订历史}}
\rhead{\Navigation}
\addcontentsline{toc}{chapter}{修订历史}

\begin{longtable}[l]{ | m{2cm} | m{1.5cm} | m{12cm} | }\hline

    \rowcolor{lightgray}
    日期 & 版本 & \hypertarget{Release notes}{发布说明}\\\hline
    \endfirsthead

    \multicolumn{3}{c}{\bfseries 接上页}\\\hline
    \rowcolor{lightgray}
    日期 & 版本 & 发布说明\\\hline
    \endhead

    \multicolumn{3}{r}{\textbf{见下页}}\\
    \endfoot

    \endlastfoot

    % table contents

    2025-xx-xx & v1.2 &
    更新以下章节：
    \begin{itemize}
    \item 章节 \ref{mod:efuse} \textit{\nameref{mod:efuse}}：更新 \ref{sec:register} \textit{\nameref{sec:register}} 和表 \ref{tab:efuse-block0-para} %DOC-10339
    \item 章节 \ref{mod:ecdsa} \textit{\nameref{mod:ecdsa}}：将模块从 ECDSA 加速器更名为 ECDSA 数字签名外设 (ECDSA\_DS)；移除 EFUSE\_ECDSA\_FORCE\_USE\_HARDWARE\_K 相关描述
    \item 章节 \ref{mod:i2c} \textit{\nameref{mod:i2c}}：更新从机启动方式的描述 %DOC-12530
    \item 章节 \ref{mod:rmt} \textit{\nameref{mod:rmt}}：更新 RMT 时钟频率限制相关公式 %DOC-13581
    \item 章节 \ref{mod:digsig} \textit{\nameref{mod:digsig}}：将模块从数字签名算法 (DSA) 更名为 RSA 数字签名外设 (RSA\_DS) %DOC-13834
    \end{itemize}
    \\\hline
    2025-07-01 & v1.1 &
    更新以下章节：
    \begin{itemize}
        \item 章节 \ref{mod:dma} \textit{\nameref{mod:dma}}：新增 GDMA\_OUTFIFO\_OVF\_CH\regindex{n}\_INT、GDMA\_OUTFIFO\_UDF\_CH\regindex{n}\_INT、GDMA\_INFIFO\_OVF\_CH\regindex{n}\_INT 和 GDMA\_INFIFO\_UDF\_CH\regindex{n}\_INT 中断的描述 %DOC-5353
        \item 章节 \ref{mod:lowpowm} \textit{\nameref{mod:lowpowm}}：增加 RTC\_EVT\_TICK 的事件描述 %DOC-10466
        \item 章节 \ref{mod:tmpsen} \textit{\nameref{mod:tmpsen}}：更新温度传感器 ETM 任务名称 %DOC-10298
    \end{itemize}
    \\\hline
    2025-03-07 & v1.0 &
    新增以下章节：
    \begin{itemize}
        \item 章节 \ref{mod:lowpowm} \textit{\nameref{mod:lowpowm}}
        \item 章节 \ref{mod:psdet} \textit{\nameref{mod:psdet}}
        \item 章节 \ref{mod:rng} \textit{\nameref{mod:rng}}
    \end{itemize}
    更新以下章节：
    \begin{itemize}
        \item 章节 \ref{mod:resclk} \textit{\nameref{mod:resclk}}：更新关于 RC\_SLOW\_CLK 时钟的描述；增加域 \hyperref[fielddesc:LPAONFORCEDOWNLOADBOOT]{LP\_AON\_FORCE\_DOWNLOAD\_BOOT} 的描述 %DOC-9389, DOC-10330
        \item 章节 \ref{mod:intmtrx} \textit{\nameref{mod:intmtrx}}：移除 CACHE\_INTR 中断源 %DOC-8488
        \item 章节 \ref{mod:aes} \textit{\nameref{mod:aes}}：新增小节 \ref{sec:aes-pseudo-round} \textit{\nameref{sec:aes-pseudo-round}} %DOC-10338
        \item 章节 \ref{mod:ecc} \nameref{mod:ecc}：将寄存器字段前缀从 ECC 更新为 ECC\_MULT；新增小节 \ref{sec:ecc-anti-attack} \textit{\nameref{sec:ecc-anti-attack}} %DOC-10131 %DOC-10336
        \item 章节 \ref{mod:ecdsa} \nameref{mod:ecdsa}：新增小节 \ref{sec:ecdsa-anti-attack} \textit{\nameref{sec:ecdsa-anti-attack}} %DOC-10336
        \item 章节 \ref{mod:extmemencr} \textit{\nameref{mod:extmemencr}}：增加关于伪轮次抗 DPA 的描述 %%DOC-10337
        \item 章节 \ref{mod:uart} \textit{\nameref{mod:uart}}：更新 FIFO 相关描述，澄清发送块和接收块的 RAM 存储空间以及 FIFO 使用方式 %DOC-8534
        \item 章节 \ref{mod:parlio} \textit{\nameref{mod:parlio}}：新增无限制长度数据传输特性 %DOC-10335
        \item 章节 \ref{mod:tmpsen} \textit{\nameref{mod:tmpsen}} 和 \ref{mod:adc} \textit{\nameref{mod:adc}}：将原章节 \textit{SAR ADC 转换器与温度传感器} 拆分为章节 \ref{mod:tmpsen} \textit{\nameref{mod:tmpsen}} 和 \ref{mod:adc} \textit{\nameref{mod:adc}}；修正几处寄存器访问类型的笔误；将 ADC 滤波公式中的 ``-0.5" 更正为 ``+0.5" %DOC-10227 %DOC-10226
    \end{itemize}
    \\\hline
    2024-06-21 & v0.5 &
    更新以下章节：
    \begin{itemize}
        \item 章节 \ref{mod:dma} \textit{\nameref{mod:dma}}：更新 suc\_eof 和 EOF 标志的相关描述 %DOC-5478
        \item 章节 \ref{mod:resclk} \textit{\nameref{mod:resclk}}：更新 \hyperref[fielddesc:LPCLKRSTSLOWCLKSEL]{LP\_CLKRST\_SLOW\_CLK\_SEL} 字段的描述
        \item 章节 \ref{mod:systimer} \textit{\nameref{mod:systimer}}：从时间补偿相关描述中删除 Deep-sleep 模式，因为仅在从 Light-sleep 唤醒后需要时间补偿 %DOC-6662
        \item 章节 \ref{mod:permctrl} \textit{\nameref{mod:permctrl}}：更新章节标题
        \item 章节 \ref{mod:uart} \textit{\nameref{mod:uart}}：更新图 \ref{fig:uart-data-frame} \textit{\nameref{fig:uart-data-frame}}，STOP 位应始终保持高电平；更新 \ref{sec:uart-wakeup} 小节描述，将“当上升沿个数大于 \hyperref[fielddesc:UARTACTIVETHRESHOLD]{UART\_ACTIVE\_THRESHOLD} 时唤醒芯片”改为“当上升沿个数大于等于 \hyperref[fielddesc:UARTACTIVETHRESHOLD]{UART\_ACTIVE\_THRESHOLD} 时唤醒芯片” %DOC-6816
        \item 章节 \ref{mod:ledpwm} \textit{\nameref{mod:ledpwm}}：更新表 \ref{tab:ledc-common-freq-res} 中的最低精度  %DOC-7304
    \end{itemize}
    \\\hline
    2023-10-12 & v0.4 &
    新增章节 \ref{mod:ecdsa} \textit{\nameref{mod:ecdsa}} \newline
    更新章节 \ref{mod:resclk} \textit{\nameref{mod:resclk}}：新增 \hyperref[fielddesc:PCRFOSCTICKNUM]{PCR\_FOSC\_TICK\_NUM} 字段的描述
    \\\hline
    2023-08-02             & v0.3 &
    新增章节 \ref{mod:iomuxgpio} \textit{\nameref{mod:iomuxgpio}} \newline
    更新以下章节：
    \begin{itemize}
        \item 章节 \ref{mod:resclk} \textit{\nameref{mod:resclk}}：在 \ref{sec:resclk-periperal-clock-configuration} 小节下方新增一条关于外设时钟门控和复位的说明 %DOC-5822
        \item 章节 \ref{mod:bootctrl} \textit{\nameref{mod:bootctrl}}：更正 \ref{sec:bootModeControl} \textit{\nameref{sec:bootModeControl}}  %DOC-5013
        \item 章节 \ref{mod:intmtrx} \textit{\nameref{mod:intmtrx}}：更新章节 \ref{sec:intmtx-int-terminology} \textit{\nameref{sec:intmtx-int-terminology}} %doc-4877
        \item 章节 \ref{mod:evntaskmatrix} \textit{\nameref{mod:evntaskmatrix}}：将支持事件任务矩阵的外设由 RTC 看门狗定时器改为 RTC 定时器 %DOC-5151
    \end{itemize}
    新增章节 \hyperref[interrupt-config-registers]{\textit{中断配置寄存器}}
    \\\hline
    2023-06-29        & v0.2 &
    新增以下章节：
    \begin{itemize}
        \item 章节 \ref{mod:riscvtracenc} \textit{\nameref{mod:riscvtracenc}}
        \item 章节 \ref{mod:resclk} \textit{\nameref{mod:resclk}}
        \item 章节 \ref{mod:ecc} \textit{\nameref{mod:ecc}}
        \item 章节 \ref{mod:i2s} \textit{\nameref{mod:i2s}}
    \end{itemize}
    更新以下章节：
    \begin{itemize}
        \item 章节 \ref{mod:riscvcpu} \textit{\nameref{mod:riscvcpu}}：更正中断 ID %DOC-5390
        \item 章节 \ref{mod:efuse} \textit{\nameref{mod:efuse}}：更新EFUSE\_SEC\_DPA\_LEVEL 寄存器的描述 %DOC-5323
        \item 章节 \ref{mod:intmtrx} \textit{\nameref{mod:intmtrx}}：删除 MSPI\_INTR 中断源及其映射寄存器，原因是 MSPI 模块暂未作说明 %doc-5368
        \item 章节 \ref{mod:permctrl} \textit{\nameref{mod:permctrl}}：在表 \ref{tab:pmp-apm} 和图 \ref{fig:pmp_apm} 中增加 EX\_MEM，从章节 \ref{sec:permctrl-features} \textit{\nameref{sec:permctrl-features}} 中删除“外部存储器”，因为 APM 没有对其的权限管理
    \end{itemize}
    新增章节 \hyperref[programming-reserved-field]{\textit{如何配置寄存器的保留域}}
    \\\hline
    2023-05-24 & v0.1 & 首次发布\\\hline

\end{longtable}
